\section{Method}
\label{sec:method}

Our model is the continuous bitstream diffusion model of \citet{method:cobit} (CoBit), a
text-generation framework, applied \emph{unchanged} to sequences that interleave several
modalities. The corruption process, the denoiser parameterisation and its loss, the noise
schedule and its estimator, the self-conditioning mechanism, the sampler and the backbone are
those of the text-only model, with the same hyperparameters. What is added is small and lives
entirely in the data pipeline: (i) a way of writing any set of tokenised modalities into one
shared binary vocabulary, so that a single network reads and writes all of them through one
interface (\S\ref{sec:method:multimodal}); and (ii) a training recipe under which one set of
weights supports unconditional co-generation of all modalities present, conditional generation
in every direction between them, and any mixture of the two, with the direction chosen at
inference time by deciding which positions are held fixed (\S\ref{sec:method:directions}).
There are no modality-specific encoders, decoders, heads, embedding tables or losses, and no
modality-specific training stages. We recap the base model first, to the extent needed to make
these two additions precise.

\subsection{Continuous bitstream diffusion}
\label{sec:method:cobit}

\textbf{Bits as the universal representation.} CoBit models a sequence of $N$ discrete tokens
$t_1,\dots,t_N$ from a vocabulary of size $V$ by writing each token as its
$m=\lceil\log_2 V\rceil$-bit binary index, giving a flat bitstream $\vx_0\in\{0,1\}^{S}$ with
$S=Nm$. Nothing else about the data is used: the framework does not know whether a token is a
word piece, an image patch code, a speech frame or a residue, and it never sees a token
identity, only its bits. Once written as bits, the sequence is treated as continuous data. A
standard variance-exploding Gaussian diffusion is run on it, $\vx_\sigma=\vx_0+\sigma\,\bm{\epsilon}$
with $\bm{\epsilon}\sim\mathcal N(\vzero,\mathbf I)$ and
$\sigma\in[\sigma_{\min},\sigma_{\max}]=[0.002,80]$, taking the data centre to be $1/2$ and
$\sigma_{\mathrm{data}}=1/2$. A denoiser $D_\theta(\vx_\sigma,\sigma)\in(0,1)^S$ predicts the
posterior probability that each bit is one and is trained with the EDM-weighted denoising loss
\citep{image:edm} in probability space,
\begin{equation}
\Ls(\theta)=\E\Big[w(\sigma)\,\tfrac1S\big\|D_\theta(\vx_\sigma,\sigma)-\vx_0\big\|^2\Big],
\qquad
w(\sigma)=\frac{\sigma^2+\sigma_{\mathrm{data}}^2}{\sigma^2\sigma_{\mathrm{data}}^2},
\label{eq:loss}
\end{equation}
which induces the score $\nabla\log p_\sigma(\vx_\sigma)\approx(D_\theta-\vx_\sigma)/\sigma^2$
used by the sampler. The network is not asked to learn the local Gaussian denoising problem:
the analytic posterior logit of an isolated bit under a uniform prior,
$(\vx_\sigma-\tfrac12)/\sigma^2$ (clipped to $\pm30$), is added to a learned contextual
residual, so that the network models only the dependence \emph{between} bits---which is where
all modality-specific structure resides. Three components of the base recipe matter for what
follows.

\textbf{Entropy-derived noise schedule.} The rate at which information about $\vx_0$ is
destroyed as $\sigma$ grows depends on the data, and different modalities---and mixtures of
modalities---lose their information at different noise levels. Rather than fixing a proposal
over $\sigma$ by hand, CoBit measures this rate on the data being modelled and allocates both
training noise levels and sampling steps according to it. The conditional entropy rate per
unit $\log\sigma$ is estimated online from the unweighted per-example denoising error,
$\hat h(\sigma)\propto e(\sigma)/\sigma^2$ with $e(\sigma)=\tfrac1S\|D_\theta-\vx_0\|^2$, kept
in a FIFO buffer over recent training examples ($8\times10^5$) and binned into 128 log-spaced
bins. The training proposal for $\sigma$, and through its inverse CDF the sampling grid, is
$q(\sigma)\propto g(\sigma;c,n)\,\hat h(\sigma)^{1/2}$, where $g(\sigma;c,n)=\sigma^n/(\sigma^n+c^n)$
($c=0.1$, $n=3$) is a smooth gate removing mass from very small $\sigma$. Training starts from
the EDM log-normal proposal ($P_{\mathrm{mean}}=-1.2$, $P_{\mathrm{std}}=1.2$), ramps linearly to
$q$ over steps $4$--$5\times10^4$, and refreshes $q$ every $2\times10^3$ steps thereafter.
Because $\hat h$ is measured rather than assumed, the schedule is self-calibrating: the same
estimator with the same hyperparameters is used for every modality combination, and it adapts
itself to each.

\textbf{Self-conditioning.} Following \citet{method:analogbits}, the denoiser takes a second
input $\hat\vx_0$, an estimate of the clean bits, concatenated bitwise with $\vx_\sigma$. During
training, with probability $p_{\mathrm{sc}}=0.5$ a first no-gradient pass produces
$\hat\vx_0=D_\theta(\vx_\sigma,\sigma,\vzero)$, which is detached and fed to the second, trained
pass; otherwise $\hat\vx_0=\vzero$. At sampling time two refresh modes are available. In
\texttt{carry} mode the probabilities returned at step $i$ are passed as $\hat\vx_0$ to step
$i{+}1$, so self-conditioning costs no additional network evaluations; in \texttt{refined} mode
an extra denoiser pass is run at the new state $(\vx_{\sigma_{i+1}},\sigma_{i+1})$ conditioned on
the carried estimate, doubling the evaluations per step. We use \texttt{carry} throughout
unless stated otherwise.

\textbf{Sampler and backbone.} We integrate the probability-flow ODE
$\mathrm d\vx/\mathrm d\sigma=(\vx-D_\theta(\vx,\sigma))/\sigma$ with a first-order
(DDIM/Euler) step \citep{method:ddim} on the entropy-derived grid, optionally with EDM-style
churn (re-noising by a factor $1+\gamma_i$ before each step; \citealp{image:edm}) as the
stochastic variant. Integration stops at a terminal $\sigma_{\mathrm{dec}}>\sigma_{\min}$; the
final probabilities are thresholded at $1/2$, regrouped into $m$-bit codes and handed to the
appropriate tokenizer's decoder. The number of steps, the guidance scale
(\S\ref{sec:method:directions}) and the churn are inference-time choices that require no
retraining. The denoiser is the Sequence Diffusion Transformer of \citet{method:cobit}: the
$m$ bits of a position, with their per-bit positional features, form one patch that is linearly
projected to one trunk token, so the transformer runs at the token length $N$ rather than the
bit length $Nm$; the trunk uses RoPE, SwiGLU and adaLN-zero conditioning on $\sigma$, and a light
skip-connection head combines the trunk output with the raw noisy bits to produce the $m$
per-bit logits of each position. The same backbone is used for every modality combination;
only its depth and width vary between checkpoints.

\subsection{Multimodal bitstreams}
\label{sec:method:multimodal}

The base model needs nothing from its input except that it be a sequence of fixed-width binary
codes. This is what makes it modality-agnostic, and it makes extending it to new modalities a
data-preparation step rather than a modelling one. Three ingredients are required.

\textbf{Modality-specific tokenizers.} Each modality is mapped to a run of integer codes by a
discrete tokenizer of its own---a byte-pair vocabulary for text, a vector- or lookup-free
quantiser for images or audio, a residue alphabet or a structure quantiser for proteins---and
back by the matching decoder. The tokenizers are the \emph{only} modality-specific components
in the system. They are frozen, they are not trained or fine-tuned jointly with the diffusion
model, and they sit outside it: the diffusion model never sees an embedding, a spectrogram or a
pixel, only bits. Any tokenizer that yields fixed-width codes can be used; quantisers whose
codebook is the full set of sign patterns of a latent vector (lookup-free or finite-scalar
quantisation; \citealp{image:lfq,audio:fsq}) pair especially naturally with bitstream
diffusion, since every bit string is then a valid code and the bits are the quantiser's own
latent coordinates rather than an arbitrary index. The tokenizers used for each modality pair
are stated in the corresponding experimental subsection (\S\ref{sec:exp}).

\textbf{One shared code space.} The per-modality vocabularies are placed in disjoint ranges of
a single integer code space, together with a small reserved block of control codes (padding and
segment markers), and $m$ is chosen as the bit width of that union---in practice by offsetting
the vocabularies or, equivalently, by prepending one or more namespace bits that select the
modality. The $m$ bits of a position therefore determine both the modality and the content of a
token, every position is a ``token of any modality'', and the patch size of the backbone is $m$
for all of them. Consequently there are no modality-specific embedding tables, encoders, heads
or losses, and the shared Bernoulli output head costs $O(m)=O(\log V)$ per position regardless
of how many vocabularies are merged. The one modality-aware parameter we add is a learned
segment embedding: a zero-initialised vector per segment type (control, text, image, audio,
\dots) added to the per-bit trunk input, which tells the network which segment a bit belongs to
without introducing any per-modality computation.

\textbf{A fixed sequence template.} A training example is laid out as a fixed-length template
of segments, one per modality, each of a fixed token budget (truncated or padded), separated by
structural markers (start/end of sequence and of each segment). Markers are ordinary codes in
the reserved block; they are always presented clean, at training and at sampling, and are never
denoised. Two-dimensional or higher-dimensional token grids are flattened along a
locality-preserving (Hilbert) curve so that 1-D adjacency reflects spatial adjacency. Modalities
that are aligned position by position (for instance a residue identity and a structure token
describing the same site) may also share one position by concatenating their codes into a single
$m$-bit word; this too is only a layout choice.\note{Confirm against the final protein codec in
\S\ref{sec:exp:protein}: the current stub describes an 18-bit position formed by a 5-bit residue
code and a 13-bit DPLM-2 structure code.} The transformer sees only the flat bitstream; the same
code path decodes each segment through its own tokenizer's decoder.

\textbf{What is, and is not, modality-specific.} Across every experiment in this paper the
corruption process, the $\sigma$ range and $\sigma_{\mathrm{data}}$, the loss and its weighting,
the entropy-schedule estimator and all of its hyperparameters, self-conditioning, the sampler
family and the backbone architecture are identical.\note{Verify against the final audio and
protein training configurations before submission; this sentence is the paper's central claim
and should be checked field by field.} What changes between modality pairs is confined to the
data pipeline: the tokenizers and their decoders, the segment layout (which modalities, in
which order, with what token budgets), and the list of training regimes of
\S\ref{sec:method:directions}. The only quantities chosen \emph{per direction} at inference are
the sampler settings (step count, guidance scale, churn), which are free hyperparameters of any
diffusion sampler and involve no retraining and no architectural difference.

\subsection{Generative directions}
\label{sec:method:directions}

A sequence built as above is partitioned into segments $\gS_1,\dots,\gS_K$ (one per modality,
plus markers). Any generation task over these modalities is a choice of a \emph{conditioning
set} $\gC$ of segments that are given and a \emph{target set} $\gT$ that must be produced:
$\gC=\emptyset$ is unconditional joint co-generation of everything; conditioning on a text
segment gives text-to-image or text-to-speech; conditioning on an image or a speech segment
gives captioning or transcription; conditioning on a prefix of a segment gives continuation; and
so on. CoBit implements every such task with the same network and the same weights through a
single mechanism, which we call \emph{partial denoising}: the positions in $\gC$ are held at
their clean values $\vx_0$ throughout, and diffusion acts only on the positions in $\gT$.
Because the model is a bidirectional denoiser over the whole sequence, a clean segment is simply
part of the input it conditions on; no task token, task-specific head or encoder is involved.

\textbf{Training.} For each training example we draw a task from a small list of
\emph{regimes}---conditioning sets with sampling probabilities---and construct the noisy input
accordingly. Let $M\in\{0,1\}^S$ be the bit mask of the conditioned positions (the segments in
$\gC$, always including the markers). The input to the denoiser is
$\vx_\sigma=M\odot\vx_0+(1-M)\odot(\vx_0+\sigma\bm{\epsilon})$, i.e.\ the conditioning segments
are clamped clean rather than noised, and the loss~\eqref{eq:loss} is averaged only over the
denoised positions $1-M$. When self-conditioning is active the clean conditioning bits are also
written into $\hat\vx_0$ on the masked positions, so the network never has to predict what it is
given. To enable classifier-free guidance \citep{method:cfg}, with probability
$p_{\mathrm{uncond}}=0.1$ the conditioning segments (but never the markers) are replaced by the
null value $1/2$---the data centre and the point of zero information for a bit under Gaussian
corruption---so that the network also learns the unconditional denoiser for each target set.
The regime list is the whole specification of which directions a model supports: adding a
direction for a given modality set means adding an entry to this list, without touching the
model. Each experimental subsection states the regimes used. Examples whose conditioning was
dropped are excluded from the entropy-rate buffer of \S\ref{sec:method:cobit}, so the schedule
reflects the conditional denoising problems the model is actually asked to solve.

\textbf{Sampling.} At inference the task fixes $M$ and the clean values on the conditioned
positions; the sampler of \S\ref{sec:method:cobit} initialises the target positions from noise,
clamps the conditioned positions to their clean bits before every denoiser call, zeroes the ODE
velocity on those positions after every call (so they never move) and clamps them once more in
the returned probabilities. Guidance uses the standard two-branch estimate in probability
space: a conditional branch sees the clean conditioning, an unconditional branch sees the null
value $1/2$ on the same positions (markers protected), the two are evaluated as one batch of
size $2B$, and the guided estimate $D_g=D_u+w\,(D_c-D_u)$ replaces $D_\theta$ in the score, where
$w=1$ recovers the plain conditional model and $w>1$ sharpens it. In \texttt{carry} mode each
branch carries its own self-conditioning estimate. Joint co-generation ($\gC=\emptyset$) uses a
single branch, so its cost per step is one network evaluation while every guided conditional
direction costs two; we report the number of function evaluations (NFE) accordingly. Since the
direction is decided entirely by $M$, the remaining sampler choices---step count, guidance
scale, churn---can be set per direction on the same weights, and nothing in the framework
requires them to agree across directions. The same weights therefore act as an unconditional
co-generator, as a family of conditional generators in every direction seen in training, and as
a model that fills in an arbitrary subset of positions given the rest.
